% Propietario conservado: /Users/ruben/Documents/New project/output/EXTREMA_MEDIA_RAZON_RECIPROCIDAD_ES_EN_20260918/ARTICULO/ES/source/libro/reciprocidad_espectral.tex
\chapter{Réciprocité thêta--Mellin et conservation de la structure polaire}
\label{x_reciprocidad:rec:espectral}

\section{Périodisation du réseau engendré}

Le réseau \(\Lambda_t\) du théorème~\ref{x_reciprocidad:rec:thm:red} possède une lecture analytique qui conserve sa dépendance au registre. Nous définissons, pour \(x>0\),
\begin{equation}
 \Theta_t(x)=\sum_{v\in\Lambda}
       \exp\!\left(-\pi_{\HMT}x\|G_tv\|^2\right).
 \label{x_reciprocidad:rec:eq:theta}
\end{equation}
Dans cette réalisation, la coordonnée \(\pi_{\HMT}\), préalablement engendrée, est reconnue avec la normalisation circulaire de Fourier. Cette reconnaissance se produit après la construction du réseau et du flux. Pour abréger l’écriture analytique, nous utiliserons \(\pi\) dans ce sens.

Sur tout intervalle compact de \(t\), les valeurs singulières de \(G_t\) sont majorées et minorées par des constantes positives. Le réseau possède un nombre \(O(R^{24})\) de points dans une boule de rayon \(R\). La comparaison par couronnes avec une gaussienne démontre la convergence absolue et uniforme de \eqref{x_reciprocidad:rec:eq:theta}, et de ses dérivées, lorsque \(x\) reste dans un compact de \((0,\infty)\). Le mode \(v=0\) contribue exactement pour un.

\begin{theorem}[Réciprocité de la lecture thêta]
\label{x_reciprocidad:rec:thm:theta}
Pour tout \(x>0\) et \(t\in\RR\),
\begin{equation}
 \Theta_t(x)=x^{-12}\Theta_{-t}(1/x).
 \label{x_reciprocidad:rec:eq:theta-reciproca}
\end{equation}
\end{theorem}
\begin{proof}
Avec la convention \(\widehat f(\xi)=\int f(y)e^{-2\pi i\langle y,\xi\rangle}\,dy\), la gaussienne \(f_x(y)=e^{-\pi x\|y\|^2}\) en dimension vingt-quatre a pour transformée \(\widehat f_x(\xi)=x^{-12}e^{-\pi\|\xi\|^2/x}\). Nous périodisons sur \(\Lambda_t\). Intégrer la périodisation sur un domaine fondamental montre que le coefficient de Fourier indexé par \(\xi\in\Lambda_t^*\) est \(\widehat f_x(\xi)/\operatorname{covol}(\Lambda_t)\). Les deux séries gaussiennes convergent absolument ; évaluer le développement à l'origine donne
\[
 \sum_{v\in\Lambda_t}f_x(v)
 =\operatorname{covol}(\Lambda_t)^{-1}
   \sum_{\xi\in\Lambda_t^*}\widehat f_x(\xi).
\]
On substitue \(\operatorname{covol}(\Lambda_t)=1\) et \(\Lambda_t^*=\Lambda_{-t}\). La démonstration conserve le mode constant dans les deux membres.
\end{proof}

La puissance \(x^{-12}\) provient du rang vingt-quatre. L’inversion du paramètre est une conséquence du réseau dual ; elle n’est pas ajoutée comme une symétrie indépendante de la lecture.

\section{Prolongement méromorphe et équation fonctionnelle}

Pour \(\Re s>12\), soit
\begin{equation}
 \mathcal E_t(s)=\sum_{v\in\Lambda\setminus\{0\}}\|G_tv\|^{-2s},
 \qquad \mathcal L_t(s)=\pi^{-s}\Gamma(s)\mathcal E_t(s).
 \label{x_reciprocidad:rec:eq:epstein}
\end{equation}
Le comptage réticulaire précédent démontre la convergence absolue dans ce demi-plan. L'intégrale gamma, appliquée terme à terme, produit
\[
 \mathcal L_t(s)=\int_0^\infty(\Theta_t(x)-1)x^{s-1}\,dx.
\]

\begin{theorem}[Décomposition polaire et réciprocité bilatérale]
\label{x_reciprocidad:rec:thm:mellin}
La fonction \(\mathcal L_t\) se prolonge méromorphiquement au moyen de
\begin{align}
 \mathcal L_t(s)={}&\frac1{s-12}-\frac1s\nonumber\\
 &+\int_1^\infty\left[(\Theta_t(x)-1)x^{s-1}
                 +(\Theta_{-t}(x)-1)x^{11-s}\right]\,dx.
 \label{x_reciprocidad:rec:eq:mellin-partida}
\end{align}
Ses seuls pôles sont simples, en \(s=0,12\), de résidus \(-1,+1\), respectivement. On a
\begin{equation}
 \mathcal L_t(s)=\mathcal L_{-t}(12-s),\qquad
 \Xi_t(s):=s(s-12)\mathcal L_t(s)=\Xi_{-t}(12-s),
 \label{x_reciprocidad:rec:eq:funcional}
\end{equation}
et \(\Xi_t\) est entière.
\end{theorem}
\begin{proof}
Nous séparons l’intégrale en \(x=1\). Pour \(0<x<1\), la réciprocité thêta donne
\[
 \Theta_t(x)-1=x^{-12}-1+x^{-12}(\Theta_{-t}(1/x)-1).
\]
Les deux premiers termes s'intègrent en \(1/(s-12)-1/s\). Dans le troisième, le changement \(y=1/x\) produit \(\int_1^\infty(\Theta_{-t}(y)-1)y^{11-s}\,dy\). La longueur minimale positive du réseau garantit la décroissance exponentielle de \(\Theta_t(x)-1\) en \(x\to\infty\). Sur chaque compact de \(s\), cette décroissance domine toute puissance de \(x\) et de \(\log x\) ; les intégrales et leurs dérivées complexes convergent donc uniformément et définissent des fonctions entières. Les résidus se lisent dans la partie rationnelle. Remplacer simultanément \((t,s)\) par \((-t,12-s)\) échange les intégrales et laisse la partie rationnelle invariante. Multiplier par \(s(s-12)\) élimine les deux pôles.
\end{proof}

Le centre \(s=6\) résulte du rang et de l’exposant \(-2s\) adopté dans \eqref{x_reciprocidad:rec:eq:epstein}. La famille est la complétion d’Epstein des réseaux déformés. Sa variable spectrale et sa normalisation restent distinctes de celles de la fonction zêta de Riemann ; la relation prouvée entre les deux constructions est la réciprocité de leurs réalisations circulaires et réticulaires.

\begin{corollary}[Conservation polaire et composante antisymétrique]
\label{x_reciprocidad:rec:cor:polar}
Les fonctions
\[
 \mathcal L_t(s)-\mathcal L_0(s),\qquad
 \mathcal A_t(s)=\frac{\mathcal L_t(s)-\mathcal L_{-t}(s)}2
\]
sont entières. De plus,
\[
 \mathcal A_t(12-s)=-\mathcal A_t(s),\qquad \mathcal A_t(6)=0.
\]
\end{corollary}
\begin{proof}
La partie rationnelle de \eqref{x_reciprocidad:rec:eq:mellin-partida} est indépendante de \(t\), de sorte qu’elle s’annule dans les deux différences. L’équation fonctionnelle change le signe de la seconde. L’évaluation en son point fixe donne \(2\mathcal A_t(6)=0\).
\end{proof}

La déformation réside donc dans une différence entière et conserve exactement les données polaires. Le zéro central appartient à la différence bilatérale, sans imposer de zéro à chaque \(\mathcal L_t\). Les pôles décrits sont des singularités d’une transformée de Mellin ; la connexion physique avec l’opérateur de masses s’établit plus loin au moyen de la section d’action, avec sa variable et son lecteur propres.

\section{Reconstruction stable à partir de l’archive bilatérale}
\label{x_reciprocidad:rec:archivo}

La loi de récupération précédente fournit, pour le décalage cyclique unitaire \(S\),
\[
 T_j=\frac{8I+S^j}{9},\qquad D_j=\frac{\sqrt8(S^j-I)}9,
 \quad j=3,4,
\]
\begin{equation}
 m=(D_3K,D_4T_3K)=MK,\qquad LM=P_{11},\qquad\|L\|=\frac94.
 \label{x_reciprocidad:rec:eq:archivo}
\end{equation}
En conséquence,
\begin{equation}
 u=P_3Lm.
 \label{x_reciprocidad:rec:eq:recuperacion-u}
\end{equation}
La carte d’incidence, le réseau marqué et le rapport régional \(\rho\) étant fixés, cette égalité retrouve successivement \(u/\|u\|\), \(g_t\), \(G_t\), \(\Lambda_t\), \(\Theta_t\) et \(\mathcal E_t\). La reconstruction de cette famille n’exige ni le terminal \(T_4T_3K\) ni la charge uniforme du registre. Les antécédents géométriques restent inclus parmi les données fixées.

\begin{theorem}[Stabilité géométrique et spectrale de la récupération]
\label{x_reciprocidad:rec:thm:estabilidad}
Supposons \(m'=m+e\), \(\|e\|\le\varepsilon\) et \(9\varepsilon/4<\|u\|\). Soit \(u'=P_3Lm'\), et reconstruisons les flux avec \(u'/\|u'\|\), sans faire varier \(\rho,P_3,L\) ni les cartes. Si \(a=|t\log\rho|\), alors
\begin{equation}
 \|g'_t-g_t\|,\ \|G'_t-G_t\|
 \le\frac92\,a e^a\frac{\varepsilon}{\|u\|}.
 \label{x_reciprocidad:rec:eq:estabilidad}
\end{equation}
Avec \(H_K=\operatorname{diag}(h_i)\), \(b=|t|\|H'_K-H_K\|\), on a
\begin{align}
 b&\le\frac{9a\varepsilon}{2\|u\|},\nonumber\\
 \Theta_t(e^{2b}x)&\le\Theta'_t(x)\le\Theta_t(e^{-2b}x),\label{x_reciprocidad:rec:eq:theta-error}\\
 e^{-2\sigma b}\mathcal E_t(\sigma)&\le\mathcal E'_t(\sigma)
             \le e^{2\sigma b}\mathcal E_t(\sigma),\qquad\sigma>12.
 \label{x_reciprocidad:rec:eq:epstein-error}
\end{align}
\end{theorem}
\begin{proof}
La norme du récupérateur donne \(\|u'-u\|\le9\varepsilon/4\), ce qui garantit \(u'\ne0\). L’inégalité triangulaire et l’inégalité inverse pour les normes impliquent
\[
 \left\|\frac{u'}{\|u'\|}-\frac{u}{\|u\|}\right\|
 \le\frac{2\|u'-u\|}{\|u\|}.
\]
Chaque exposant est compris entre \(-a\) et \(a\). Appliquer le théorème des accroissements finis à l'exponentielle et prendre la norme diagonale prouve \eqref{x_reciprocidad:rec:eq:estabilidad} ; répéter chaque entrée deux fois ne change pas cette norme. La même estimation avant exponentiation donne la borne de \(b\).

Comme les deux générateurs sont diagonaux dans la même carte,
\[
 e^{-2b}\|G_tv\|^2\le\|G'_tv\|^2\le e^{2b}\|G_tv\|^2.
\]
Appliquer l’exponentielle décroissante et sommer prouve \eqref{x_reciprocidad:rec:eq:theta-error}. Élever à \(-\sigma\) et sommer dans le domaine de convergence prouve \eqref{x_reciprocidad:rec:eq:epstein-error}. Les bornes s’appliquent à tous les vecteurs, et non à une troncature finie des séries.
\end{proof}

L’erreur \(\varepsilon\) se mesure dans les mémoires normalisées de \eqref{x_reciprocidad:rec:eq:archivo}. Si les registres rationnels \(z=m/\sqrt8\) sont conservés, leur erreur est multipliée par \(\sqrt8\) avant d’appliquer l’estimation. Les erreurs indépendantes dans le rapport régional ou dans les cartes requièrent leurs propres bornes. Le résultat récupère une famille de géométries et de lectures analytiques, avec un contrôle quantitatif de l’archive qui la détermine.

